\documentclass[11pt,a4paper]{article}
\usepackage[margin=1in]{geometry}
\usepackage{hyperref}
\usepackage{enumitem}
\usepackage{graphicx}
\usepackage{xcolor}

% -----------------------------------------------------
% bvraghav-tiet-ucs503p-article.sty
% -----------------------------------------------------

% Variable: Lab Instructor
\makeatletter \newcommand{\BvrSetLabInstructor}[1]{%
  \gdef\Bvr@LabInstructor{#1}}
% default
\newcommand{\Bvr@LabInstructor}{Dr. Nisha Thakur}
% public accessor
\newcommand{\BvrLabInstructorValue}{\Bvr@LabInstructor}
\makeatother

% Add subtitle
\usepackage{titling}
% -- Set title bold
\pretitle{\begin{center}\bfseries\fontsize{18bp}{18bp}\selectfont}
\makeatletter
\newcommand{\BvrSubtitle}[1]{%
  \posttitle{%
    \par\end{center}
    \begin{center}\large#1\end{center}
    \vskip 0.5em}%
}
\makeatother

% Hints in the section title(s)
\newcommand{\BvrHint}[1]{%
  {\normalfont\normalsize\itshape\hfill #1}}

% -----------------------------------------------------

\title{ItiNova: An AI-Assisted Personalized Travel Planning System}

\BvrSetLabInstructor{Dr. Nisha Thakur}

\BvrSubtitle{%
  Project Proposal for UCS503P  \\
  Submitted to: \BvrLabInstructorValue \\ \bfseries
  Thapar Institute of Engineering and Technology%
}

\author{
  Shreya, [CSED]%
  \thanks{Roll No: \texttt{1024030520}, %
    Email: \texttt{sshreya1\_be24\@thapar.edu}}
  \and
  Molika Goela, [CSED]%
  \thanks{Roll No: \texttt{1024030524}, %
    Email: \texttt{mgoela\_be24\@thapar.edu}}
  \and
  Jiya Gupta, [CSED]%
  \thanks{Roll No: \texttt{1024030501}, %
    Email: \texttt{jgupta1\_be24\@thapar.edu}}
  \and
  Himani Nayyar, [CSED]%
  \thanks{Roll No: \texttt{1024030510}, %
    Email: \texttt{hnayyar\_be24\@thapar.edu}}
}

\date{\today}

\begin{document}
\maketitle

\section{Higher-order goal \BvrHint{\ldots secondary framing}}

This project aims to reduce the time and effort required to plan a personalized trip by providing users with a unified travel-planning platform. ItiNova allows users to either select a destination they already know or discover suitable destinations based on their preferences such as mood, travel style, budget, duration, starting location, and maximum travel distance. The system then helps users explore relevant places and activities, build and customize their own itinerary, and receive AI-assisted modifications along with weather-aware packing and destination-specific health information.

The immediate engineering objective is to provide a practical, personalized, and user-controlled travel-planning workflow that combines conventional web technologies, external travel information services, and AI assistance.

\section{Time-to-value \BvrHint{\ldots primary heuristic}}

\begin{itemize}
\item We will first implement the core destination selection and trip-planning workflow so that users can create a trip and build an itinerary from relevant place and activity suggestions.
\item We will prioritize the core user experience before implementing supporting features such as packing assistance, health information, and map visualization.
\item Weather will be integrated as contextual information for itinerary planning and packing rather than as a separate standalone feature.
\item AI will be introduced as an assistance layer for itinerary modification and personalization rather than making the system fully dependent on AI-generated plans.
\item Development will follow incremental iterations so that a functional trip-planning workflow is available early and can be tested before additional modules are integrated.
\end{itemize}

\section{Problem Statement}

Travellers often need to search across multiple sources to decide where to travel, identify suitable places and activities, and manually organize them into a practical itinerary. This becomes more difficult when users have specific preferences regarding budget, travel distance, mood, interests, and travel pace.

Common pain points include:

\begin{itemize}
\item \textbf{Destination uncertainty:} users may know their preferred budget, travel distance, mood, or environment but may not know which destination would suit them.
\item \textbf{Generic recommendations:} available travel recommendations may not sufficiently account for individual interests, budget, trip duration, travel style, and preferred pace.
\item \textbf{Manual itinerary planning:} travellers often need to search for attractions and activities separately and manually decide how to arrange them across different days.
\item \textbf{Dynamic weather conditions:} weather can affect outdoor activities and what travellers need to carry, but this information is often considered separately from trip planning.
\item \textbf{Packing uncertainty:} generic packing lists do not account for the destination, trip duration, expected weather, and planned activities.
\item \textbf{Health and disease awareness:} travellers may not be aware of destination-specific health risks, seasonal disease concerns, and basic precautions before travelling.
\item \textbf{Travel-distance management:} selecting places without considering their geographical locations can result in inefficient itineraries and unnecessary travel.
\item \textbf{Frequent itinerary changes:} users may want to modify their plans according to changing preferences without rebuilding the entire itinerary manually.
\end{itemize}

\textbf{Why this matters now (contextual relevance):} Travel planning involves information from multiple sources and requires users to make several decisions before arriving at a practical plan. A personalized system can reduce repetitive searching while keeping the user in control of the final itinerary.

\section{Proposed Solution \BvrHint{\ldots Software Proposition}}

\subsection{Overview}

Build a web-based ``ItiNova'' system with the following components:

\begin{itemize}

\item \textbf{Destination discovery:} users who already know where they want to travel can directly enter their destination. Users who are uncertain can instead provide preferences such as starting location, budget, duration, maximum travel distance, mood, preferred environment, interests, and travel style. The system recommends destinations that match these preferences.

\item \textbf{Place and activity suggestions:} after selecting a destination, the system retrieves relevant attractions, activities, and places that match the user's interests and trip requirements. These suggestions act as the building blocks for the user's itinerary.

\item \textbf{Personalized itinerary planning:} users can select suggested places and activities and create their own day-by-day itinerary. They can add, remove, edit, and reorder activities according to their preferences.

\item \textbf{AI-assisted itinerary modification:} an AI module assists users in modifying their itinerary through natural-language requests. Users can request changes such as making a day less hectic, adding more adventure activities, rearranging activities, or adapting the plan to their preferences.

\item \textbf{Weather-aware planning:} weather information for the selected travel period is used as contextual information while planning activities and generating packing recommendations. Weather is not treated as a separate standalone feature.

\item \textbf{Smart packing assistant:} the system generates a personalized packing checklist using destination, trip duration, expected weather, and planned activities. Users can mark items as packed and modify the checklist.

\item \textbf{Travel health and disease awareness:} the system provides general destination-specific health information, relevant disease or health-risk awareness, precautions, and a basic pre-travel checklist. The feature is intended for general travel guidance and is not a medical diagnosis system.

\item \textbf{Interactive travel map:} the map displays the locations included in the user's itinerary and allows users to visualize the geographical arrangement of planned activities and routes.

\item \textbf{Trip saving and history:} users can save their created trips and access previously planned itineraries for future reference.

\end{itemize}

\subsection{Core workflow}

\begin{enumerate}

\item The user enters the application and chooses either:
\begin{itemize}
    \item a known destination, or
    \item the ``Help Me Choose'' option when they are unsure where to travel.
\end{itemize}

\item If the user chooses ``Help Me Choose'', they provide information such as starting location, budget, duration, maximum travel distance, mood or vibe, preferred environment, interests, and travel style.

\item The system processes these preferences and recommends suitable destinations based on the user's major constraints and preferences.

\item The user selects a destination and provides or confirms trip details such as travel dates, duration, number of travellers, budget, interests, and preferred travel pace.

\item The system retrieves relevant places and activities for the selected destination using external travel information services.

\item Relevant weather and geographical information are retrieved for the selected trip period and destination.

\item The user views the suggested places and activities and creates their itinerary by selecting, adding, removing, editing, and reordering activities according to their needs.

\item The AI module can assist the user in modifying the created itinerary through natural-language instructions such as making a particular day less hectic, adding more activities of a particular type, or rearranging activities according to weather and travel preferences.

\item The packing assistant generates a personalized checklist using destination, trip duration, expected weather, and planned activities.

\item The system provides destination-specific general health and disease awareness information together with relevant precautions and a basic travel health checklist.

\item The itinerary locations are displayed on an interactive map so that users can visualize their planned activities and travel routes.

\item The completed trip and itinerary can be saved so that the user can access and modify it later.

\end{enumerate}

\subsection{Operational constraints for an educational travel-planning setting}

\begin{itemize}

\item Keep the solution usable on low-to-mid range devices through responsive web design.

\item Avoid maintaining a large manually curated travel database. Relevant places, geographical information, and weather data will be retrieved through external services when required.

\item Use an existing large language model/API rather than training a new foundation model.

\item Keep external API integrations modular so that failure of an individual service does not affect the entire application.

\item Keep the user in control of the itinerary. AI-generated modifications will be treated as suggestions that users can accept, reject, or further edit.

\item Present health and disease information as general travel guidance and not as medical diagnosis or treatment recommendations.

\item Use current external information for dynamic travel information such as weather, locations, and routes instead of relying solely on the language model's stored knowledge.

\end{itemize}

\section{Solution Approach \BvrHint{\ldots Engineering focus}}

This is an engineering course project focused on building a practical AI-assisted travel-planning workflow rather than developing a new machine learning model. The system will combine conventional web development, external travel information services, a database for persistent application data, and an existing large language model.

\subsection{Destination discovery and place suggestions \BvrHint{\ldots personalized search}}

The destination discovery component will support two different user scenarios.

\begin{itemize}

\item \textbf{Known destination:} users can directly enter the destination they want to visit.

\item \textbf{Unknown destination:} users can provide preferences such as starting location, budget, duration, maximum travel distance, mood, environment, interests, and travel style. The system will use these preferences to recommend suitable destinations.

\end{itemize}

After the destination is selected, relevant places and activities will be retrieved through external travel information services rather than maintaining a large manually curated database.

The retrieved information may include:

\begin{itemize}
\item Place name and location
\item Place category or type
\item Ratings where available
\item Geographical coordinates
\item Relevant activities
\item Distance and travel-time information
\end{itemize}

These suggestions will then be displayed to the user as possible building blocks for their itinerary.

\subsection{User-controlled itinerary planning \BvrHint{\ldots practical personalization}}

The initial itinerary will be created and controlled by the user rather than being automatically generated by AI.

Users will be able to:

\begin{itemize}
\item Select suggested places and activities.
\item Add activities to particular days.
\item Remove activities.
\item Edit activity details.
\item Reorder activities.
\item Move activities between different days.
\item Adjust the itinerary according to their own preferences.
\end{itemize}

The system will use available geographical and weather information to provide useful context while the user plans the trip.

\subsection{AI-assisted itinerary modification \BvrHint{\ldots practical AI integration}}

The AI component will use an existing large language model/API rather than training a new foundation model.

The AI will not be responsible for making the user's initial itinerary. Instead, it will assist users when they want to modify or improve their existing plan.

The backend can provide the AI with:

\begin{itemize}
\item Existing itinerary
\item Destination
\item Travel dates and duration
\item Budget
\item Number of travellers
\item Interests
\item Mood or travel vibe
\item Preferred travel pace
\item Relevant places and activities
\item Weather information
\item Travel distance and route information
\end{itemize}

\subsection{Weather-aware planning \BvrHint{\ldots contextual information}}

Weather will not be implemented as an independent travel recommendation module.

Instead, weather information will be used as contextual information for:

\begin{itemize}
\item Selecting suitable activities for particular days.
\item Helping users decide between outdoor and indoor activities.
\item Assisting AI-based itinerary modifications.
\item Generating a personalized packing checklist.
\end{itemize}

Current or forecast weather information will be retrieved through an external weather service rather than relying on the language model's internal knowledge.

\subsection{Smart packing assistant \BvrHint{\ldots contextual assistance}}

The packing assistant will generate a personalized checklist using:

\begin{itemize}
\item Destination
\item Trip duration
\item Weather conditions
\item Planned activities
\item General travel requirements
\end{itemize}

The checklist will be displayed in an interactive format where users can mark items as packed, remove unnecessary items, and add their own items.

\subsection{Travel health and disease awareness \BvrHint{\ldots responsible guidance}}

The health module will provide general destination-specific health and disease awareness information.

The system may consider:

\begin{itemize}
\item Destination
\item Travel season
\item Relevant environmental conditions
\item General destination-specific health risks
\item Basic precautions
\end{itemize}

The feature will be presented as general travel guidance and will not provide medical diagnosis or treatment recommendations.

\subsection{Interactive travel map \BvrHint{\ldots geographical context}}

The interactive map will display locations included in the user's itinerary.

The map will allow users to:

\begin{itemize}
\item View itinerary locations.
\item Understand the geographical arrangement of activities.
\item View routes between planned locations.
\item Explore the destination visually.
\end{itemize}

\subsection{Web architecture \BvrHint{\ldots web-based solution}}

A typical 3-tier web architecture will be used:

\begin{itemize}

\item \textbf{Frontend:} responsive interface for destination selection, preference collection, place suggestions, itinerary creation and editing, packing checklist, health information, map visualization, and saved trips.

\item \textbf{Backend API:} authentication, trip management, destination recommendation logic, external API communication, AI integration, itinerary modification, weather integration, and data validation.

\item \textbf{Database:} store users, trip information, preferences, created itineraries, and saved trip data. The database will not be used as a large repository of all possible destinations and attractions.

\end{itemize}

External services will be accessed through the backend so that API credentials are not directly exposed to the frontend.

\section{Evaluation Criterion \BvrHint{\ldots measurable and attributable}}

We define success using measurable metrics related to the core travel-planning workflow.

\subsection{Primary evaluation metric}

\textbf{Time-to-itinerary (TTI):} time between the user submitting their trip requirements and obtaining a usable itinerary containing selected activities.

\begin{itemize}
\item \textbf{Target:} median TTI $\leq$ 60 seconds during the pilot window.
\item \textbf{Attribution:} measured using application timestamps between trip submission and the availability of the initial usable itinerary.
\end{itemize}

\subsection{Secondary evaluation metrics}

\begin{itemize}

\item \textbf{Destination recommendation relevance:} user rating of recommended destinations on a 1--5 scale.

\item \textbf{Place suggestion relevance:} user rating of the relevance of suggested places and activities on a 1--5 scale.

\item \textbf{Itinerary usefulness:} user rating of the usefulness of the created itinerary on a 1--5 scale.

\item \textbf{User control:} percentage of users who successfully add, remove, edit, or reorder activities without assistance.

\item \textbf{AI modification usefulness:} user rating of AI-assisted itinerary modifications on a 1--5 scale.

\item \textbf{Packing relevance:} percentage of generated packing items that users consider relevant.

\item \textbf{Health-information usefulness:} user rating of the usefulness and clarity of the provided health and disease awareness information.

\item \textbf{System reliability:} successful completion rate of core workflows such as destination selection, place retrieval, itinerary creation, itinerary modification, and trip saving.

\end{itemize}

\subsection{Pilot validation plan \BvrHint{\ldots high-level}}

\begin{itemize}

\item Recruit 10--20 pilot users as student travellers.

\item Ask users to create sample trips using different destinations, budgets, durations, and travel preferences.

\item Measure the time required to create an initial usable itinerary.

\item Collect 1--5 ratings for destination recommendation relevance, place suggestion relevance, itinerary usefulness, AI modification usefulness, packing relevance, and overall planning experience.

\item Ask users to perform common itinerary modifications such as adding, removing, editing, and reordering activities.

\item Use application logs to measure response times and workflow completion.

\end{itemize}

\section{Scalability \BvrHint{\ldots with foresight}}

The system should remain practical as the number of users and saved trips increases.

\begin{enumerate}

\item \textbf{Scalable (theoretically):} the system should support growth through:
\begin{itemize}
\item separation of frontend and backend components,
\item indexed database queries,
\item stateless API design where appropriate,
\item modular external API integrations,
\item selective retrieval of external travel information,
\item storing only required user and trip data instead of maintaining a large manually curated travel database.
\end{itemize}

\item \textbf{Operationally deployable (practically):} the system should be deployable using:
\begin{itemize}
\item a managed database,
\item environment-based configuration for API credentials,
\item a web-hosted frontend and backend,
\item modular external API services.
\end{itemize}

\end{enumerate}

\section{Engine availability heuristic}

A mature set of tools and services is available to implement the proposed system:

\begin{itemize}

\item Web framework for responsive frontend development and REST APIs.

\item Node.js and Express for backend development.

\item MongoDB for persistent application data.

\item Existing large language model/API for itinerary modification and personalization.

\item Current weather API for forecast information.

\item Current places and geocoding API for destination and place information.

\item Routing API for geographical distance and travel-time information.

\item Open-source map technologies for itinerary visualization.

\item Authentication and session management for user accounts.

\item Testing tools for backend and application functionality.

\end{itemize}

These mature components allow the project to focus on workflow design, practical AI assistance, personalized travel planning, usability, and system correctness rather than developing new infrastructure or training a foundation AI model.

\section{Project scope and deliverables \BvrHint{\ldots iteration-friendly}}

\subsection{Initial deliverable \BvrHint{\ldots first iteration}}

\begin{itemize}

\item User registration/login.

\item Trip creation form.

\item Known-destination trip planning.

\item ``Help Me Choose'' destination discovery based on starting location, budget, distance, duration, mood, interests, and travel style.

\item Destination recommendation and selection.

\item Retrieval of relevant places and activities.

\item User-controlled itinerary creation.

\item Basic itinerary editing, deletion, and reordering.

\item Weather information integration as contextual trip information.

\item Saving generated trips and itineraries.

\end{itemize}

\subsection{Subsequent deliverables}

\begin{itemize}

\item AI-assisted itinerary modification through natural-language requests.

\item Smart packing assistant based on destination, duration, activities, and weather.

\item Interactive travel map showing itinerary locations and routes.

\item Travel health and disease awareness information.

\item Basic travel health and safety checklist.

\item Improved destination and place recommendations.

\item Improved AI itinerary modification using weather and geographical information.

\item Improved validation of itinerary constraints.

\end{itemize}

\section{Risks and mitigations}

\begin{itemize}

\item \textbf{External API availability:} weather, places, geocoding, or routing services may have rate limits or temporary failures. Use modular service layers, limited API calls, and fallback handling.

\item \textbf{AI-generated information:} AI may generate inaccurate or unsuitable itinerary modifications. Present AI output as a recommendation, validate structured responses, and allow users to review and edit all changes.

\item \textbf{AI response latency:} AI modification requests may take time. Display an appropriate loading state and measure response time during evaluation.

\item \textbf{Outdated AI knowledge:} dynamic travel information should not be obtained solely from the language model. Current weather, place, and route information will be retrieved from external services and provided as context to the AI.

\item \textbf{Recommendation quality:} destination or place suggestions may not always match user expectations. Allow users to modify their preferences and manually select or remove suggested activities.

\item \textbf{API quota limitations:} free external APIs may impose daily request limits. API calls will be minimized by retrieving only relevant information and avoiding unnecessary repeated requests.

\item \textbf{Health information accuracy:} health and disease information can be sensitive and time-dependent. The system will present it as general awareness information and avoid medical diagnosis or treatment recommendations.

\item \textbf{Data privacy:} only information required for account management and trip planning will be stored.

\end{itemize}

\section{Summary}

This project proposes \textbf{ItiNova}, an AI-assisted personalized travel-planning platform designed to reduce the time and effort required to plan a trip.

The system supports two primary entry points: users can either provide a destination they already know or receive destination recommendations based on preferences such as starting location, budget, travel distance, duration, mood, interests, and travel style.

After selecting a destination, the system retrieves relevant places, activities, weather, and geographical information. Users can then create their own itinerary by selecting, adding, removing, editing, and reordering activities according to their preferences.

AI is used as an assistance layer rather than as the sole decision-maker. Users can provide natural-language instructions to modify their existing itinerary, such as making a day less hectic, adding more adventure activities, or rearranging activities based on weather and preferences.

Supporting features include a smart packing assistant, destination-specific health and disease awareness, and an interactive travel map. Weather is treated as contextual information that supports itinerary planning and packing rather than as a separate standalone feature.

The project focuses on practical AI integration, external travel information services, user-controlled personalization, and a deployable web-based architecture without requiring a large manually maintained travel database or development of a new AI model.

\end{document}